\subsection{代码数据的特殊地位}

代码数据在预训练中有着超出”编程能力”的作用。Folklore 认为，代码训练有助于模型的\textbf{推理能力}——因为代码本身就是结构化的逻辑表达。GitHub 代码的处理有几个特殊挑战：

\begin{itemize}
    \item \textbf{重复率极高}：fork、复制粘贴、自动生成的代码（如 protobuf、配置文件）使得原始 GitHub 数据中有大量近重复。The Stack 项目从 1.37 亿个仓库中克隆了 510 亿个文件，但去重后只有 50 亿个是唯一的。
    \item \textbf{许可证过滤}：The Stack 只保留宽松许可（MIT、Apache）的代码，使用 go-license-detector 工具自动识别。这将数据量从 TB 级压缩到 3.1TB。
    \item \textbf{不只是代码}：仓库中还有 README、issue、commit message、pull request 评论等元数据，这些对训练也有价值。GH Archive 以小时为单位记录 GitHub 事件快照。
\end{itemize}

